\documentclass[12pt, a4paper]{article}

% PACKAGES
\usepackage[utf8]{inputenc}
\usepackage{amsmath}
\usepackage{amssymb}
\usepackage[margin=1.2in]{geometry}
\usepackage{authblk}
\usepackage{bm}
\linespread{1.2}

\title{A Reflection on the Magnetic Field: The Microscopic Texture of the Computational Grid}
\author{Haifeng Qiu}
%\date{}

\begin{document}

\begin{center}
    {\uppercase{\Large \textbf{A Reflection on the Magnetic Field}}} \\
    \vspace{0.5em}
    {\large \textit{The Microscopic Texture of the Computational Grid}} \\
    \vspace{2em}
    \textbf{Haifeng Qiu} \\
    \today
    \vspace{2em}
\end{center}



\begin{abstract}
This brief reflection clarifies the existential form of the magnetic field within the framework of ``Computational Realism.'' We propose that while the electric field manifests as the macroscopic gradient of the computational field, the magnetic field emerges as its dynamic microscopic texture—specifically, a Moiré pattern. The paper consists of a single conjecture that reconstructs the origin of magnetic force through the lens of non-equilibrium dynamics.
\end{abstract}



\vspace{2em}


\section{The Core Conjecture: ``Gliding'' Over an Un-equilibrated Grid}

Within the discrete vacuum grid of the computational field, the \textbf{electric field} manifests as the macroscopic gradient, while the \textbf{magnetic field} emerges as its dynamic microscopic texture—an \textbf{``unbalanced texture of entropy momentum''} carved jointly by the finite propagation speed of information (retarded potentials) and the kinematic distortion (Lorentz effects and Doppler effect). It represents the residual, imperfectly equilibrated micro-details of a dynamic system, and the so-called \textbf{magnetic force} is ultimately nothing more than the macroscopic manifestation of a moving charge \textbf{``gliding through''} and structurally undergoing orthogonal deflection as it interacts with these grid textures.


\section{A Conceptual Proof: Multi-body Dynamic Equilibrium and the Inevitability of the ``Workless'' Residue}

In classical electrodynamics, the force exerted by a magnetic field on a charge (the Lorentz force, $\bm{F} = q\bm{v} \times \bm{B}$) is always perpendicular to the direction of the particle's motion and does no work on the particle. This section attempts to explore the physical origin of this specific mathematical form within the framework of a conceptual proof. Due to the complexity of microscopic details, a complete and rigorous mathematical derivation cannot be provided here. Instead, the aim is to illustrate how, under the constraints of a finite speed of light and the underlying information grid, this ``orthogonal and workless'' form of force emerges as the inevitable residue after the dynamic equilibration of a multi-body system.

\subsection{Symmetry Breaking of the Moiré Pattern: Time Integral and Path Integral}

In a discrete computational grid where the speed of information propagation is limited to the speed of light $c$, the microscopic momentum flows (electric fields) emitted by positive and negative source charges in different states of motion cannot achieve perfect cancellation in a vacuum. This microscopic cancellation deficit, caused by retardation, weaves into a high-frequency, periodic wave texture in spacetime, which we heuristically understand as a \textit{``phase Moiré pattern.''}

\begin{itemize}
    \item \textbf{The Immunity of a Stationary Charge:} A stationary test charge situated within this microscopic texture constantly experiences the impact of microscopic momentum flows. However, because the Moiré pattern possesses strict time periodicity, the \textbf{time integral} of these alternating positive and negative phase shifts \textbf{periodically returns to zero}.
    Therefore, while a stationary charge oscillates at the microscopic level, it experiences no net force on a macroscopic scale.

    \item \textbf{The Force on a Moving Charge:} When a test charge cuts into this non-linear grid texture formed by the Doppler effect with a macroscopic velocity $\bm{v}$, its own motion breaks the temporal symmetry of its experience. As its motion path undergoes non-linear distortion while slicing through these microscopic phases, the essence of the force it experiences transforms into a path integral along that trajectory. This effect, analogous to the \textbf{Ponderomotive force} mechanism, results in a non-zero accumulation of spatial phase differences that no longer vanishes in a state of motion, thereby manifesting as a real macroscopic force. This explains the fundamental mechanism underlying the magnetic force and why it is only effective on moving charges.
\end{itemize}

\subsection{The Filtering of Interaction Forms by Multi-body Self-Equilibration Mechanisms}

A moving charge experiences a force due to the path integral, and the reason this force must inevitably manifest in an orthogonal and workless form is constrained by the multi-body self-equilibration mechanisms within the entity generating the macroscopic current (such as a wire or plasma).

Here, we introduce a \textbf{basic dynamical assumption}: within the propagation and response range permitted by the speed of light, any unbalanced momentum that can be responded to and cancelled out by the motion of the vast number of free charges within the system will be rapidly absorbed and flattened internally, thus preventing it from leaking significantly into the macroscopic external space. Based on this assumption, we dynamically filter the possible interaction forms of the ``microscopic unbalanced momentum'' radiated outward by the source:

\begin{enumerate}
    \item \textbf{Cancellation of Static Components:} Suppose this microscopic Moiré pattern contains components capable of exerting a macroscopic thrust on a stationary charge. As this component propagates outward, the free charges inside the wire will immediately feel this thrust and undergo microscopic displacements. This minuscule redistribution of charge will generate an internal electric field in the opposite direction, balancing out the portion of electrostatic momentum attempting to escape. Therefore, the macroscopic field that ultimately dissipates cannot exert a force on static charges.
    
    \item \textbf{Cancellation of Work-Doing Components:} Suppose the microscopic Moiré pattern contains components capable of doing parallel work (accelerating or decelerating) on moving charges. When it passes through other parallel or anti-parallel electron flows inside the wire, it will inevitably do work on these electrons and change their kinetic energy. According to the conservation of energy, after being subjected to a longitudinal force, the electron flow will immediately excite an opposing ``negative-work feedback field'' (macroscopically analogous to the effect of self-inductance). This feedback field, established at the speed of light, will absorb and flatten all fluctuations attempting to output energy longitudinally. Thus, any microscopic momentum residue attempting to do work will be consumed by the dynamic currents within the system itself. Furthermore, since the charges within the wire lack any macroscopic motion in the vertical direction, they cannot exert a work-doing influence on external charges moving in that direction.  In fact, because neither the positive ions nor the electrons possess macroscopic vertical velocity, the vertical force components they exert on a vertically moving test charge perfectly cancel each other out, resulting in a net vertical force of zero.
\end{enumerate}

\subsection{Conclusion: The Residue That Does No Work on a Single Moving Charge}

After passing through the two filtering mechanisms formed by multi-body interactions inside the wire---aimed at static forces and dynamic work-doing components---the only specific interaction form left in the residual momentum of the microscopic grid that can penetrate the charge flow within the wire and dissipate into the vacuum is this: \textbf{It can only exert an effect on moving charges, and it does no work on a single moving charge.}

Geometrically, doing no work on a single particle means that the direction of the force must always be strictly perpendicular to the particle's velocity ($\bm{F} \perp \bm{v}$). Such a force can neither directly change the kinetic energy of the particles nor be completely cancelled out by the electron flow restricted to one-dimensional motion inside the wire.

Within the exploratory picture of \textit{``Computational Realism,''} the magnetic force described in classical electromagnetism can be understood as the workless residue that inevitably leaks into macroscopic space after the multi-body moving charge system on a finite-light-speed grid has completed all possible internal work and equilibration. This incompletely equilibrated microscopic texture constitutes the physical entity of the magnetic field.



\section{A Simplified Two-Body Model: Kinematic Estimation of Magnetic Field Strength for Parallel and Anti-Parallel Moving Charges}

In the previous section, through the neutralization and self-equilibration mechanisms of macroscopic multi-body systems, we demonstrated the orthogonal geometric form of the magnetic field as a ``workless residue.'' However, the complete solution of multi-body dynamical equations is extremely complex, making it difficult to directly derive the absolute strength of the magnetic field.

To quantitatively estimate the magnitude of this ``residual interaction,'' it is necessary to temporarily strip away the macroscopic background and adopt an extremely simplified vacuum two-body model. We attempt to demonstrate that the ``magnetic field strength'' in classical electromagnetism can be completely derived directly from purely kinematic mechanisms on the underlying computational grid, such as \textbf{information propagation delay (retardation)}, \textbf{spatial projection}, \textbf{bit-change flux enhancement}, and the \textbf{Doppler effect}.

\subsection{Core Commonalities: Diagonal Retardation of Entropy Momentum and Lorentz Kinematics}

Whether moving co-directionally or anti-parallel to each other, when two isolated charges move at a speed $v$ on the computational grid and are at a certain instant in the same transverse plane (with a separation distance of $L$), their interaction is governed by the following two fundamental dynamical commonalities of the grid:

\textbf{First, the diagonal propagation and geometric attenuation of microscopic entropy momentum.} \\
Due to the absolute speed limit $c$ of information propagation on the grid, the microscopic entropy momentum (electric field information) perceived by the receiver at this exact moment must have been emitted by the emitter $\Delta t$ time ago. During this retardation time, the actual trajectory traveled by the signal on the grid is a diagonal line $r$. \\
According to the geometric relationship $r^2 = L^2 + (v \Delta t)^2$ and $r = c \Delta t$, the actual propagation distance is elongated to:
\[
r = \frac{L}{\sqrt{1 - v^2/c^2}} = \gamma L
\]
The interaction force weakens inversely with the square of the spatial distance due to diffusion, thus introducing an attenuation of $1/r^2 = 1/(\gamma^2 L^2)$. Assuming the angle between the propagation path and the direction of motion is $\theta$, the transverse projection factor is $\sin\theta = L/r = 1/\gamma$. Combining these two points, \textbf{the transverse force attenuation factor arising from pure geometry is $1/\gamma^3$}.

\textbf{Second, the dynamical effects induced by Lorentz kinematics.} \\
When a charge moves on the computational grid, its structural adaptation manifests through three distinct Lorentz effects that influence its interaction:
\begin{itemize}
    \item \textbf{Bit-change flux enhancement:} According to the definition of ``Computational Realism,'' the mass/strength of a charge is its \textbf{bit-change flux} as an ``information soliton.'' To maintain structural self-consistency while in motion, the internal computational activity of the charge must increase. This elevates the frequency of its underlying entropy waves, which enhances the interaction strength of the entropy momentum at both the emitting and receiving ends each by a Lorentz factor $\gamma$.
    \item \textbf{Length contraction:} The compression of the particle along the direction of motion geometrically concentrates its field, enhancing the emission strength in the direction perpendicular to the motion by a factor of $\gamma$. Conversely, its narrowed spatial profile weakens its reception capability by a factor of $\gamma$. When two charges move at the same speed, these morphological effects perfectly cancel each other out.
    \item \textbf{Time dilation:} The slowing down of the particle's overall macroscopic cyclic rhythm serves only to maintain internal synchronization. It does not affect the actual strength of its electric field interactions with the external environment.
\end{itemize}

The above two points form the foundation of any two-body kinematic interaction. Next, we will introduce the Doppler effect to estimate specific scenarios based on different relative directions of motion.

\subsection{Scenario One: Parallel Co-directional Motion (The Essence of Magnetic Attraction)}

Consider two charges moving parallel and co-directionally with the same velocity $v$. According to the geometric relationship in Section 3.1, the cosine of the angle between the hypotenuse of signal propagation and the longitudinal axis of motion is $\cos\theta = (v\Delta t)/(c\Delta t) = \beta$. Because the emitter is ``chasing'' the signal it emitted, the wavefront density is compressed by a factor of $1/(1-\beta\cos\theta) = 1/(1-\beta^2)$; simultaneously, because the receiver is ``fleeing'' from the incoming signal at the same speed and angle, its interception rate is reduced by a factor of $(1-\beta\cos\theta) = (1-\beta^2)$. These two competing effects on the grid perfectly cancel each other out. Therefore, the Doppler effect factor in this system is $1$.

By multiplying all the above factors, we obtain the transverse interaction force $F_{y(parallel)}$ during parallel co-directional motion:
\[
F_{y(parallel)} = F_{static} \cdot \underbrace{\left(\frac{1}{\gamma^3}\right)}_{\text{Geometry}} \cdot \underbrace{(\gamma \cdot \gamma)}_{\text{Flux}} \cdot \underbrace{(\gamma \cdot 1/\gamma)}_{\text{Morphology}} \cdot \underbrace{(\frac{1-\beta^2}{1-\beta^2} )}_{\text{Doppler}}
\]
Which simplifies to:
\[
F_{y(parallel)} = \frac{F_{static}}{\gamma}
\]
This result clearly indicates: the transverse repulsive force between parallel, co-directionally moving charges must necessarily be weaker than the Coulomb repulsive force $F_{static}$ in a state of absolute rest. In the phenomenological description of classical electromagnetism, this \textbf{relative weakening} of the repulsive force is interpreted as the generation of an additional ``magnetic attractive term.''

\subsection{Scenario Two: Anti-Parallel Relative Motion (The Surge of Magnetic Repulsion)}

Now, consider two charges moving anti-parallel to each other with velocities $+v$ and $-v$ respectively, and let us calculate the interaction force at the instant they cross paths (with a transverse separation of $L$).

In this case, the Doppler effects no longer cancel out but produce an intense non-linear superposition. According to the geometric relationship in Section 3.1, the cosine of the angle between the hypotenuse of signal propagation and the longitudinal axis of motion is $\cos\theta = (v\Delta t)/(c\Delta t) = \beta$. Because the emitter is ``chasing'' the signal it emitted, the wavefront density is compressed by a factor of $1/(1-\beta\cos\theta) = 1/(1-\beta^2)$; simultaneously, the receiver is ``colliding'' head-on with the wavefront at a relative velocity, increasing its interception rate by a factor of $(1+\beta\cos\theta) = (1+\beta^2)$. \\
Therefore, the \textbf{total Doppler enhancement factor} at the instant of anti-parallel motion is:
\[
\frac{1+\beta^2}{1-\beta^2} = \gamma^2(1+\beta^2)
\]
Substituting this Doppler factor into the overall equation, we calculate the transverse interaction force $F_{y(anti-parallel)}$ at the instant of intersection:
\[
F_{y(anti-parallel)} = F_{static} \cdot \underbrace{\left(\frac{1}{\gamma^3}\right)}_{\text{Geometry}} \cdot \underbrace{(\gamma \cdot \gamma)}_{\text{Flux}} \cdot \underbrace{(\gamma \cdot 1/\gamma)}_{\text{Morphology}} \cdot \underbrace{\left[ \gamma^2(1+\beta^2) \right]}_{\text{Doppler}}
\]
Which simplifies to:
\[
F_{y(anti-parallel)} = F_{static} \cdot \gamma(1+\beta^2)
\]
This result demonstrates: the transverse repulsive force between anti-parallel moving charges will be \textbf{significantly greater} than the Coulomb repulsive force in a stationary state. In macroscopic classical observation, this directly corresponds to the phenomenon where ``anti-parallel currents repel each other.''

\subsection{Scenario Three: Moving Charge Acting on a Stationary Charge (Relativistic Electric Field Enhancement)}

Consider a source charge moving at velocity $v$ and a test charge at rest, at the instant they are separated by a transverse distance $L$. In this asymmetric state of motion, the four kinematic factors on the grid do not cancel out, but instead lead to a net intensification of the interaction.

First, as the emitter is moving, its signal follows the diagonal retardation path described in Section 3.1, introducing a geometric attenuation factor of $1/\gamma^3$. Second, the emitter’s internal bit-change flux is enhanced by $\gamma$, while the stationary receiver remains at unity. Third, the morphology of the moving emitter causes a lateral ``squeezing'' of its field lines, increasing the field density by a factor of $\gamma$. Finally, because the emitter is ``chasing'' its own signal at the moment of emission ($\cos\theta = \beta$), the Doppler wavefront compression factor is $1/(1-\beta^2) = \gamma^2$, while the stationary receiver has a neutral interception factor of 1. \\
Multiplying these grid factors together, we obtain the force exerted by the moving charge on the stationary one:
\[
F_{m \to s} = F_{static} \cdot \underbrace{\left(\frac{1}{\gamma^3}\right)}_{\text{Geometry}} \cdot \underbrace{(\gamma \cdot 1)}_{\text{Flux}} \cdot \underbrace{(\gamma \cdot 1)}_{\text{Morphology}} \cdot \underbrace{(\gamma^2 \cdot 1)}_{\text{Doppler}}
\]
Which simplifies to:
\[
F_{m \to s} = \gamma F_{static}
\]
This result is consistent with the Lorentz transformation of the electric field in classical electrodynamics, where the transverse field of a moving point charge is enhanced by $\gamma$.

\subsection{Scenario Four: Stationary Charge Acting on a Moving Charge (The Invariance of the Static Field)}

Now, let us reverse the roles: consider a stationary source charge acting on a test charge moving at velocity $v$. At the instant of their closest approach $L$, we analyze the interaction from the perspective of the grid's absolute rest frame.

In this case, the source is at rest; thus, the signal propagates along a direct transverse path with no diagonal retardation ($\text{Geometry} = 1$) and no Doppler compression ($\text{Doppler} = 1$). However, the moving receiver undergoes its own structural adaptations. Its internal bit-change flux is enhanced by a factor of $\gamma$ to maintain self-consistency, but its spatial profile along the direction of motion is simultaneously contracted, narrowing its ``interception cross-section'' for incoming transverse signals by a factor of $1/\gamma$. \\
The total interaction force is the product of these competing endogenous effects:
\[
F_{s \to m} = F_{static} \cdot \underbrace{(1)}_{\text{Geometry}} \cdot \underbrace{(1 \cdot \gamma)}_{\text{Flux}} \cdot \underbrace{(1 \cdot 1/\gamma)}_{\text{Morphology}} \cdot \underbrace{(1 \cdot 1)}_{\text{Doppler}}
\]
Which simplifies to:
\[
F_{s \to m} = F_{static}
\]
This concludes that a stationary charge exerts a force on a moving charge that is \textbf{identical} to the Coulomb force at rest.

\subsection{Conclusion}

Through the estimation of the two basic kinematic modes above, without introducing any independent ``magnetic field vector,'' we have successfully derived the magnitude scale and directional characteristics of classical magnetic phenomena. The strength of the magnetic field is ultimately the result of a joint interplay of the underlying electric field (microscopic entropy momentum) on a finite-light-speed grid.

In this framework, the motion of an individual charge establishes a quasi-continuous, distorted electric field grid that inherently encapsulates both electric and magnetic interactions. In contrast, the ``pure'' magnetic field produced by a macroscopic current is the inevitable residue after multi-body self-equilibration, manifesting as a periodic, anisotropic texture at the charge-lattice scale. When a test charge moves across this grid texture, it's instantaneous spatiotemporal coordinate and its own directional Lorentz and Doppler distortions interact with the grid’s anisotropy. This geometric complementarity facilitates a precise, orthogonal coupling, resulting in a ``workless'' force perpendicular to the direction of motion. On a macroscopic scale, this underlying grid texture can be effectively modeled as the magnetic vector potential ($\mathbf{A}$), while the force arising from this complementary interaction is precisely identified as the magnetic force. 



\section{The Emergence of ``Relativity'' under an Absolute Grid}

To verify whether this theory satisfies the principle of relativity within an \textbf{absolute spatiotemporal grid}, we integrate the four kinematic factors of the system into a universal formula and analyze the transformation laws of internal interaction forces when the entire physical system translates at a global velocity $V$.

\subsection{Establishing the Universal Force Formula in the ``Absolute Grid''}

On the absolutely stationary computational grid, assume the absolute velocity of the emitting charge (source) is $v_1$, and the absolute velocity of the receiving charge (test charge) is $v_2$ (assuming both move horizontally). Based on the previous analysis, the total transverse interaction force $F_{grid}$ is the product of the following four factors: \textbf{geometric retardation} ($1/\gamma_{v1}^3$), \textbf{flux enhancement} ($\gamma_{v1} \cdot \gamma_{v2}$), \textbf{morphological contraction} ($\gamma_{v1} / \gamma_{v2}$), and the \textbf{Doppler effect} ($\gamma_{v1}^2 (1 - \beta_{v1}\beta_{v2})$).

Integrating these factors, the universal absolute force formula is:
\[
F_{grid} = F_{static} \cdot \underbrace{\left[ \frac{1}{\gamma_{v1}^3} \right]}_{\text{Geometry}} \cdot \underbrace{\left[ \gamma_{v1} \cdot \gamma_{v2} \right]}_{\text{Flux}} \cdot \underbrace{\left[ \frac{\gamma_{v1}}{\gamma_{v2}} \right]}_{\text{Morphology}} \cdot \underbrace{\left[ \gamma_{v1}^2 (1 - \beta_{v1}\beta_{v2}) \right]}_{\text{Doppler}}
\]
After algebraic cancellation, the simplified form is:
\[
F_{grid} = F_{static} \cdot \gamma_{v1} (1 - \beta_{v1}\beta_{v2})
\]

\subsection{Analysis of System Behavior under Global Translation (Velocity $V$)}

Assume a physical system (such as a laboratory) translates as a whole at velocity $V$ relative to the absolute grid.

\textbf{1. Velocity Composition Logic} \\
Within the scope of studying the relativistic invariance of particle motion, the velocity of a particle refers to the observed value measured by the deformed rulers and clocks within the local reference frame. Since the rulers and clocks inside the system—as information solitons—have already undergone substantial Lorentz effects while moving at velocity $V$, we must apply algebraic corrections to convert the local observed velocity $u$ into the actual absolute velocity $v$ on the grid. This correction derives a velocity addition formula consistent with relativity:
\[
v_1 = \frac{u_1 + V}{1 + \frac{u_1V}{c^2}}, \quad v_2 = \frac{u_2 + V}{1 + \frac{u_2V}{c^2}}
\]

\textbf{2. Calculating the Actual Grid Force $F'$} \\
We need to calculate the actual force $F'$ transmitted by the underlying grid while the system is in a state of translation at velocity $V$.

\textbf{Step One: Calculate the Doppler term $(1 - \beta_{v1}\beta_{v2})$} \\
Substituting the velocity formulas and expanding:
\[
1 - \beta_{v1}\beta_{v2} = 1 - \frac{(\beta_{u1} + \beta_V)(\beta_{u2} + \beta_V)}{(1 + \beta_{u1}\beta_V)(1 + \beta_{u2}\beta_V)}
\]
Using the algebraic identity $1 - \frac{(a+x)(b+x)}{(1+ax)(1+bx)} = \frac{(1-ab)(1-x^2)}{(1+ax)(1+bx)}$, we simplify it to:
\[
1 - \beta_{v1}\beta_{v2} = \frac{1 - \beta_{u1}\beta_{u2}}{\gamma_V^2 (1 + \beta_{u1}\beta_V)(1 + \beta_{u2}\beta_V)}
\]

\textbf{Step Two: Calculate the absolute Lorentz factor $\gamma_{v1}$ at the emitter} \\
Based on the definition $\gamma_{v1} = 1/\sqrt{1-\beta_{v1}^2}$, substituting the absolute velocity $v_1$ and expanding gives:
\[
\gamma_{v1} = \gamma_{u1} \gamma_V (1 + \beta_{u1}\beta_V)
\]

\textbf{Step Three: Assemble the final grid force $F'$} \\
Substituting the above two terms into the universal grid force formula $F' = F_{static} \cdot \gamma_{v1} (1 - \beta_{v1}\beta_{v2})$:
\[
F' = F_{static} \cdot \left[ \gamma_{u1} \gamma_V (1 + \beta_{u1}\beta_V) \right] \cdot \left[ \frac{1 - \beta_{u1}\beta_{u2}}{\gamma_V^2 (1 + \beta_{u1}\beta_V)(1 + \beta_{u2}\beta_V)} \right]
\]
Canceling common factors $(1 + \beta_{u1}\beta_V)$ and one $\gamma_V$ in the numerator and denominator, we obtain:
\[
F' = \frac{F_{static} \gamma_{u1} (1 - \beta_{u1}\beta_{u2})}{\gamma_V (1 + \beta_{u2}\beta_V)}
\]

\subsection{Comparison and Physical Analysis}

Observe the numerator of the above result: $F_{static} \gamma_{u1} (1 - \beta_{u1}\beta_{u2})$. In mathematical form, this term is identical to the interaction force measured by an internal observer when the system is at \textbf{absolute rest} ($V=0$), which we denote as $F_{lab}$.

Replacing the numerator with $F_{lab}$, the equation becomes:
\[
F' = \frac{F_{lab}}{\gamma_V (1 + \beta_{u2}\beta_V)}
\]
Rearranging terms gives the transformation relationship between the force measured inside the laboratory and the true force on the absolute grid:
\[
F_{lab} = F' \cdot \gamma_V (1 + \beta_{u2}\beta_V)
\]
This mathematical form matches exactly with the transformation formula for transverse forces in \textbf{classical electrodynamics}.

\subsection{Conclusion}

The above derivation demonstrates that Lorentz covariance emerges naturally from algebraic operations involving signal propagation delay and particle structural deformation on an absolutely stationary Euclidean grid. Although the system's global translation velocity $V$ alters the particle's absolute mechanical parameters on the grid, the physical deformations of measuring instruments (\textbf{time dilation and length contraction}) and the \textbf{Doppler effect} of signals undergo precise algebraic coupling. This coupling perfectly offsets the impact of absolute motion for any internal observer. Therefore, within the framework of ``Computational Realism,'' the existence of an absolute reference frame is logically consistent with the principle of relativity; relativistic effects manifest as the inevitable macroscopic emergence of underlying discrete dynamical mechanisms.




\section{Non-Equilibrium Dynamics—Induction, Radiation, and the Geometric Localization of Photons}

\subsection{The Order Stratification of Multi-body Interactions and the System's ``Self-Equilibration'' Mechanism}

\textbf{Problem Statement:} In the aforementioned two-particle interaction model, we primarily discussed the transverse force perpendicular to the direction of motion. However, according to the derivation, a moving charge should exert a thrust on a parallel moving charge in the longitudinal direction of motion. This might seem to conflict with the observational fact in classical electromagnetism that ``steady currents do not generate longitudinal forces.'' To analyze this phenomenon, we need to perform a stratified dynamical analysis of the ``self-equilibration filtering'' mechanism within the dynamic equilibrium framework of multi-body systems, categorized by the kinematic order of the charge's motion.

\begin{enumerate}
    \item \textbf{Zero-Order Effect (Position $x$): Electrostatic Equilibrium.}\\
    The ``information fountain'' generated by a static charge manifests as an omnidirectionally diffusing gradient field. In multi-body systems such as wires, any longitudinal thrust component attempting to escape will immediately trigger a microscopic displacement of free charges (electrostatic induction), thereby being perfectly cancelled out internally within the system.

    \item \textbf{First-Order Effect (Velocity $v$): Texturization of the Steady Magnetic Field.}\\
    Constant currents interweave into microscopic ``Moiré patterns'' on the grid. Although this motion generates a residue of momentum flow, the multi-body system establishes a stable equilibrium internally through ``self-induction'' due to the constant velocity.
    \begin{itemize}
        \item \textbf{Filtering Mechanism:} All dynamical components attempting to do work (cause acceleration) on external charges are digested and flattened by the dynamic currents within the system. Specifically, if the kinematic distortion field of one moving charge has the capacity to do work on another moving charge, this portion of the field effect will inevitably be absorbed first by the moving charges internal to the system. When the system is in a steady state, charges at different positions within the wire move at a constant velocity, and the work-doing effects caused by this motion are macroscopically cancelled out due to symmetry.
        \item \textbf{Residual Product:} What ultimately penetrates the multi-body system and dissipates into the external macroscopic space is solely the \textbf{``orthogonal, workless residue''} that can neither push static charges nor change the kinetic energy of moving charges. This constitutes the mechanical origin of the classical magnetic force (Lorentz force) where $\bm{F} \perp \bm{v}$ and does no work.
    \end{itemize}

    \item \textbf{Second-Order Effect (Acceleration $a$): Equilibrium Lag and the Essence of Induction.}\\
    When a charge enters an accelerating state (second-order motion), the pre-existing multi-body self-equilibration mechanism undergoes an instantaneous collapse due to the speed limitations of information propagation and charge response, leading to an overflow of energy.
    \begin{itemize}
        \item \textbf{Cancellation Delay and Phase Lag:} The system requires time to establish a negative feedback equilibrium. When the current changes over time, velocity differences emerge among charges at different positions within the wire. The compensatory response of internal free charges is constrained by their own \textbf{dynamical inertia} and spatial asymmetry, making instantaneous cancellation impossible and thus producing a phase lag.
        \item \textbf{Overflow of Momentum Flow:} This portion of the momentum flow, which fails to be cancelled in time and possesses the capacity to do longitudinal work, breaks through the confines of the multi-body system and dissipates into the vacuum grid as an ``unbalanced residue.''
        \item \textbf{Reconstruction of Faraday's Law:} In multi-body scenarios, the macroscopically observed ``induced electric field'' is essentially the \textbf{information overflow of the accelerating charge's field caused by the system's response delay}. It is no longer bound by the ``orthogonal and workless'' constraint and can directly exert longitudinal thrust and transfer energy to other charges in space.
    \end{itemize}
\end{enumerate}

\subsection{Analysis of the Microscopic Morphology and Mechanism of Electromagnetic Wave Propagation}

\begin{enumerate}
    \item \textbf{The Physical Morphology of Electromagnetic Waves:}\\
    The physical entity of a photon is the periodic variation of the entropy momentum texture (i.e., the shape of the Moiré pattern) along its propagation path, which originates from the variation in the source charge's velocity. If the direction and intensity of this texture are represented by the magnetic vector potential, then the electromagnetic wave is the propagation of the periodic fluctuation of the magnetic vector potential in the direction perpendicular to the motion.

    \item \textbf{The Transversality of Radiation:}\\
    The momentum flow possesses a unidirectional constraint; it cannot achieve backflow or longitudinal back-and-forth density oscillations along the axis of the charge's motion. Therefore, any ``information perturbation'' excited by velocity changes cannot form a wave along the longitudinal axis. Instead, this (periodic) variation in the shape of the Moiré pattern---namely, the variation in the magnetic vector potential---propagates outward within the \textbf{transverse plane perpendicular to the flow direction} through a diffusion process similar to friction. This dynamical constraint directly leads to the transverse polarization attribute of electromagnetic waves. Furthermore, the propagation and confining effect of this varying magnetic vector potential might well be the origin of matter's particle-like nature.

    \item \textbf{The Propagation and Diffusion of Electromagnetic Waves:}\\
    Because the momentum fluctuation of a single photon is constrained to propagate within a transverse plane, its amplitude attenuates as $1/r$ with distance. This is the microscopic reason why photons can propagate over long distances in space. In contrast, the macroscopic photon stream emitted by a light source undergoes a spherical $1/r^2$ energy diffusion in three-dimensional space.

    \item \textbf{The Unification of Wave-Particle Attributes:}\\
    Photons are not completely independent entities but rather the wave-like propagation of momentum flow on the grid. Within a finite space, the fluctuations of multiple photons will overlap, diffuse, and interfere, exhibiting complete wave characteristics. Although the propagation is wave-like, its absorption process is constrained by the structural transition energy levels of the detector (atomic bound states), manifesting as a holistic, quantized absorption.
\end{enumerate}

\subsection{Conclusion of this Chapter}

We argue that the dynamical laws of classical electromagnetism are the inevitable products of the computational grid responding to different kinematic orders, under the combined constraints of the \textbf{limit of information propagation (speed of light retardation)}, \textbf{charge response delay (inertia)}, and the \textbf{directional constraint of momentum flow}.

The macroscopic \textbf{magnetic field} is the microscopic ``texture'' left after the system achieves self-equilibration; \textbf{induction} is the momentum ``overflow'' following the system's equilibrium lag; and the \textbf{photon} is the self-sustaining wave propagation effect of the geometric morphology of this overflowing momentum flow across the transverse dimension.



\end{document}